% ============================================================
% 第5章补充：经典例题与自学元素
% 插入位置：第5章现有miniquiz之后（\end{miniquiz} 之后）
% ============================================================

\section*{习题精选与详解}
\addcontentsline{toc}{section}{习题精选}

\begin{example}{托卡马克与实验室等离子体的比压计算}{beta-calc}
比压 $\beta = 2\mu_0 p / B^2$ 是衡量热压与磁压相对重要性的无量纲参数。请分别计算：
\begin{enumerate}[label=(\arabic*)]
\item 托卡马克装置：$B=5\,\mathrm{T}$，$n=10^{20}\,\mathrm{m}^{-3}$，$T_e=T_i=10\,\mathrm{keV}$；
\item 实验室低温等离子体：$B=0.1\,\mathrm{T}$，$n=10^{18}\,\mathrm{m}^{-3}$，$T_e=T_i=4\,\mathrm{eV}$。
\end{enumerate}
并解释 $\beta\ll 1$ 和 $\beta\sim 1$ 的物理意义。
\tcblower
\textbf{解答：}

\textbf{(1) 托卡马克：}
等离子体总压强 $p = p_e + p_i = 2nk_{\!B}T$（电子与离子各贡献一份）：
\[
p = 2\times 10^{20}\times 1.602\times 10^{-15} = 3.204\times 10^5\,\mathrm{Pa}
\]
磁压 $B^2/2\mu_0$：
\[
\frac{B^2}{2\mu_0} = \frac{25}{2\times 4\pi\times 10^{-7}} = \frac{25}{2.513\times 10^{-6}} \approx 9.95\times 10^6\,\mathrm{Pa}
\]
比压：
\[
\beta = \frac{2\mu_0 p}{B^2} = \frac{2\times 4\pi\times 10^{-7}\times 3.204\times 10^5}{25} = \frac{8.05\times 10^{-1}}{25} \approx 3.2\times 10^{-2}
\]
即 $\beta \approx 3.2\%$。

\textbf{(2) 实验室低温等离子体：}
\[
p = 2\times 10^{18}\times 1.602\times 10^{-19}\times 4 = 1.28\,\mathrm{Pa}
\]
\[
\beta = \frac{2\times 4\pi\times 10^{-7}\times 1.28}{0.01} = \frac{3.22\times 10^{-6}}{0.01} = 3.2\times 10^{-4}
\]

\textbf{物理意义：}
\begin{itemize}[leftmargin=*,itemsep=0pt]
\item $\beta\ll 1$：磁压远大于热压，磁场主导等离子体动力学。托卡马克中 $\beta\sim 3\%$ 已属较高，磁约束装置的工程极限通常在 $\beta<10\%$。
\item $\beta\sim 1$：热压与磁压相当，等离子体对磁场有显著的"反作用"，磁场被等离子体逆磁电流排挤出中心区域。
\item $\beta\gg 1$：热压主导，磁场几乎不影响等离子体行为（如某些天体等离子体喷流）。
\end{itemize}
\end{example}

\begin{example}{阿尔芬速度计算与对比}{alfven-speed-calc}
阿尔芬速度 $v_A = B/\sqrt{\mu_0\rho_m}$ 是磁化等离子体中信息传播的特征速度。分别计算：
\begin{enumerate}[label=(\arabic*)]
\item 托卡马克等离子体：$B=5\,\mathrm{T}$，$n=10^{20}\,\mathrm{m}^{-3}$，氘离子 $m_i=2m_p$；
\item 实验室低温氢等离子体：$B=0.1\,\mathrm{T}$，$n=10^{18}\,\mathrm{m}^{-3}$；
\item 太阳日冕：$B=10^{-3}\,\mathrm{T}$，$n=10^{14}\,\mathrm{m}^{-3}$（质子）。
\end{enumerate}
对比 $v_A$ 与等离子体中的声速 $c_s$，并说明为何阿尔芬波成为太阳日冕加热的重要候选机制。
\tcblower
\textbf{解答：}

质量密度 $\rho_m = n_i m_i$。

\textbf{(1) 托卡马克：}
\[
\rho_m = 10^{20}\times 3.344\times 10^{-27} = 3.344\times 10^{-7}\,\mathrm{kg/m^3}
\]
\[
v_A = \frac{5}{\sqrt{4\pi\times 10^{-7}\times 3.344\times 10^{-7}}} = \frac{5}{\sqrt{4.198\times 10^{-13}}} = \frac{5}{6.48\times 10^{-7}} \approx 7.7\times 10^6\,\mathrm{m/s}
\]
声速（$T=10\,\mathrm{keV}$，氘）：$c_s = \sqrt{2k_{\!B}T/m_i} = \sqrt{2\times 1.602\times 10^{-15}/3.344\times 10^{-27}} \approx 9.8\times 10^5\,\mathrm{m/s}$。
$v_A / c_s \approx 7.9$，阿尔芬速度远大于声速，磁场主导。

\textbf{(2) 实验室等离子体：}
\[
\rho_m = 10^{18}\times 1.673\times 10^{-27} = 1.673\times 10^{-9}\,\mathrm{kg/m^3}
\]
\[
v_A = \frac{0.1}{\sqrt{4\pi\times 10^{-7}\times 1.673\times 10^{-9}}} = \frac{0.1}{4.58\times 10^{-8}} \approx 2.2\times 10^6\,\mathrm{m/s}
\]
声速（$T_e=4\,\mathrm{eV}$，$T_i\approx 0.1\,\mathrm{eV}$）：$c_s = \sqrt{k_{\!B}T_e/m_i} \approx 2.0\times 10^4\,\mathrm{m/s}$。
$v_A / c_s \approx 110$，磁场仍主导。

\textbf{(3) 太阳日冕：}
\[
\rho_m = 10^{14}\times 1.673\times 10^{-27} = 1.673\times 10^{-13}\,\mathrm{kg/m^3}
\]
\[
v_A = \frac{10^{-3}}{\sqrt{4\pi\times 10^{-7}\times 1.673\times 10^{-13}}} = \frac{10^{-3}}{4.58\times 10^{-10}} \approx 2.2\times 10^6\,\mathrm{m/s}
\]
声速（$T\sim 10^6\,\mathrm{K}\sim 86\,\mathrm{eV}$）：$c_s = \sqrt{\gamma k_{\!B}T/m_i} \approx 1.5\times 10^5\,\mathrm{m/s}$。

\textbf{结论：} 在三种环境中均有 $v_A \gg c_s$，说明磁场在等离子体动力学中占绝对主导地位。太阳日冕中阿尔芬波速度高达 $2200\,\mathrm{km/s}$，且阿尔芬波可携带大量能量从光球层向上输运，是日冕加热（百万度高温）最有希望的能量来源之一。
\end{example}

\begin{example}{磁雷诺数与磁力线冻结的数值判断}{magnetic-reynolds}
磁雷诺数 $R_m = \mu_0 L v / \eta$ 衡量磁扩散与磁对流输运的相对重要性。计算以下两种情形的 $R_m$，并判断磁力线冻结是否成立：
\begin{enumerate}[label=(\arabic*)]
\item 托卡马克：$L=1\,\mathrm{m}$，$v=10^4\,\mathrm{m/s}$，$T_e=10\,\mathrm{keV}$；
\item 实验室氩气放电：$L=0.1\,\mathrm{m}$，$v=10^3\,\mathrm{m/s}$，$T_e=4\,\mathrm{eV}$。
\end{enumerate}
Spitzer电阻率近似公式：$\eta \approx 5\times 10^{-5} Z\ln\Lambda / T_e^{3/2}\,[\Omega\mathrm{m}]$，$T_e$ 单位为 $\mathrm{eV}$，取 $\ln\Lambda=10$。

\textbf{模型的适用性提示：} Spitzer 公式\emph{只计入电子--离子库仑碰撞}。下面第 (2) 问的"氩气放电"题面没有给电离度与中性气体密度，因此\textbf{这里只研究一个假设库仑碰撞主导的假想氩等离子体}；真实低温放电中电子--中性碰撞常使 $\eta$ 高出若干个数量级，绝不能把下面的 $R_m$ 直接当成实际装置的数值。这是"选模型"这一步的意义所在，请对照正文\S5.1 的"MHD 适用条件：三道独立的门槛"。
\tcblower
\textbf{解答：}

\textbf{(1) 托卡马克：}
\[
\eta \approx \frac{5\times 10^{-5}\times 10}{10^6} = 5\times 10^{-10}\,\Omega\mathrm{m}
\]
\[
R_m = \frac{\mu_0 v L}{\eta} = \frac{4\pi\times 10^{-7}\times 10^4\times 1}{5\times 10^{-10}} = \frac{1.257\times 10^{-2}}{5\times 10^{-10}} \approx 2.5\times 10^7
\]
$R_m \gg 1$，磁扩散时间 $\tau_\eta = \mu_0 L^2/\eta = R_m L/v \approx 2500\,\mathrm{s}$，远大于对流时间 $\tau_{\text{流}} = L/v = 10^{-4}\,\mathrm{s}$。磁力线冻结完全成立。

\textbf{(2) 实验室放电：}
\[
\eta \approx \frac{5\times 10^{-5}\times 10}{8} = 6.25\times 10^{-5}\,\Omega\mathrm{m}
\]
\[
R_m = \frac{4\pi\times 10^{-7}\times 10^3\times 0.1}{6.25\times 10^{-5}} = \frac{1.257\times 10^{-4}}{6.25\times 10^{-5}} \approx 2.0
\]
$R_m \sim 1$（略大于 1），磁扩散时间 $\tau_\eta = \mu_0 L^2/\eta \approx 0.2\,\mathrm{ms}$，对流时间 $\tau_{\text{流}} = L/v = 0.1\,\mathrm{ms}$。两者相当，磁力线冻结不再严格成立，磁场与流体的运动仅部分耦合——这正是实验室放电中容易观察到磁重联与磁扩散效应的原因。

\textbf{判据总结：}
\begin{itemize}[leftmargin=*,itemsep=0pt]
\item $R_m \gg 1$：磁对流主导，磁通冻结（理想MHD，如天体等离子体、托卡马克）
\item $R_m \ll 1$：磁扩散主导，磁力线可滑过流体（如实验室低参数等离子体、固态导体）
\item $R_m \sim 1$：两者竞争，出现磁重联等有趣现象（如太阳耀斑电流片）
\end{itemize}
\end{example}

\begin{example}{理想MHD中的磁能与动能转换}{mhd-energy-conversion}
从理想MHD方程组出发，证明：
\begin{enumerate}[label=(\arabic*)]
\item 磁能演化满足 $\displaystyle\frac{\partial}{\partial t}\left(\frac{B^2}{2\mu_0}\right) = -\divg\vect{S} - \vect{u}\cdot(\vect{j}\times\vect{B})$，其中 $\vect{S} = \vect{E}\times\vect{B}/\mu_0$ 为坡印廷矢量；
\item 动能演化满足 $\displaystyle\frac{\partial}{\partial t}\left(\frac{1}{2}\rho_m u^2\right) + \divg\left(\frac{1}{2}\rho_m u^2\vect{u}\right) = -\vect{u}\cdot\grad p + \vect{u}\cdot(\vect{j}\times\vect{B})$；
\item 解释 $\vect{u}\cdot(\vect{j}\times\vect{B})$ 的符号的物理意义。
\end{enumerate}
\tcblower
\textbf{解答：}

\textbf{(1) 磁能方程：}
从理想MHD磁场演化方程 $\partial\vect{B}/\partial t = \curl(\vect{u}\times\vect{B})$ 出发：
\[
\frac{\partial}{\partial t}\left(\frac{B^2}{2\mu_0}\right) = \frac{1}{\mu_0}\vect{B}\cdot\frac{\partial\vect{B}}{\partial t} = \frac{1}{\mu_0}\vect{B}\cdot\curl(\vect{u}\times\vect{B})
\]
利用矢量恒等式 $\vect{B}\cdot\curl\vect{A} = \divg(\vect{A}\times\vect{B}) + \vect{A}\cdot\curl\vect{B}$，令 $\vect{A} = \vect{u}\times\vect{B}$：
\[
\vect{B}\cdot\curl(\vect{u}\times\vect{B}) = \divg\bigl((\vect{u}\times\vect{B})\times\vect{B}\bigr) + (\vect{u}\times\vect{B})\cdot\curl\vect{B}
\]

第一项：用 BAC--CAB 公式 $\vect{a}\times(\vect{b}\times\vect{c})=\vect{b}(\vect{a}\cdot\vect{c})-\vect{c}(\vect{a}\cdot\vect{b})$，先把 $\vect{B}\times(\vect{u}\times\vect{B})$ 展开：
\[
\vect{B}\times(\vect{u}\times\vect{B}) = \vect{u}(\vect{B}\cdot\vect{B}) - \vect{B}(\vect{B}\cdot\vect{u}) = B^2\vect{u} - (\vect{u}\cdot\vect{B})\vect{B}
\]
再由反交换律 $\vect{a}\times\vect{b}=-\vect{b}\times\vect{a}$：
\[
(\vect{u}\times\vect{B})\times\vect{B} = -\vect{B}\times(\vect{u}\times\vect{B})
= -B^2\vect{u} + (\vect{u}\cdot\vect{B})\vect{B}
\]

\textbf{关键：} 理想MHD中 $\vect{E} = -\vect{u}\times\vect{B}$，取与 $\vect{B}$ 的叉乘：
\[
\vect{E}\times\vect{B} = -(\vect{u}\times\vect{B})\times\vect{B}
= B^2\vect{u} - (\vect{u}\cdot\vect{B})\vect{B}
\]
即 $\vect{E}\times\vect{B} = \mathbf{-}(\vect{u}\times\vect{B})\times\vect{B}$——\textbf{两者相差一个负号}，这正是本步最容易写错的地方。验证方式：取 $\vect{u}\perp\vect{B}$，则 $\vect{u}\times\vect{B}$ 垂直于两者，再叉乘 $\vect{B}$ 后方向回到 $\vect{u}$ 的\emph{反}方向，与 $\vect{E}\times\vect{B}$ 同向，大小 $B^2u$，即 $\vect{S}=\vect{E}\times\vect{B}/\mu_0 = B^2\vect{u}/\mu_0$。因此：
\[
\divg\bigl((\vect{u}\times\vect{B})\times\vect{B}\bigr) = -\divg(\vect{E}\times\vect{B}) = -\mu_0\divg\vect{S}
\]

第二项：$\curl\vect{B} = \mu_0\vect{j}$，且 $(\vect{u}\times\vect{B})\cdot\vect{j} = \vect{u}\cdot(\vect{B}\times\vect{j}) = -\vect{u}\cdot(\vect{j}\times\vect{B})$，故：
\[
(\vect{u}\times\vect{B})\cdot\curl\vect{B} = -\mu_0\,\vect{u}\cdot(\vect{j}\times\vect{B})
\]

综合得：
\[
\frac{\partial}{\partial t}\left(\frac{B^2}{2\mu_0}\right) = -\divg\vect{S} - \vect{u}\cdot(\vect{j}\times\vect{B})
\]

\textbf{(2) 动能方程：}
从理想MHD动量方程 $\rho_m\left(\partial\vect{u}/\partial t + \vect{u}\cdot\grad\vect{u}\right) = -\grad p + \vect{j}\times\vect{B}$，点乘 $\vect{u}$：
\[
\vect{u}\cdot\left(\rho_m\frac{\partial\vect{u}}{\partial t}\right) + \rho_m\vect{u}\cdot(\vect{u}\cdot\grad\vect{u}) = -\vect{u}\cdot\grad p + \vect{u}\cdot(\vect{j}\times\vect{B})
\]
利用 $\vect{u}\cdot\partial\vect{u}/\partial t = \frac{1}{2}\partial u^2/\partial t$ 和连续性方程，左端可化为：
\[
\frac{\partial}{\partial t}\left(\frac{1}{2}\rho_m u^2\right) + \divg\left(\frac{1}{2}\rho_m u^2\vect{u}\right)
\]
故：
\[
\frac{\partial}{\partial t}\left(\frac{1}{2}\rho_m u^2\right) + \divg\left(\frac{1}{2}\rho_m u^2\vect{u}\right) = -\vect{u}\cdot\grad p + \vect{u}\cdot(\vect{j}\times\vect{B})
\]

\textbf{(3) 物理意义：}
将 (1) 与 (2) 相加，$\vect{u}\cdot(\vect{j}\times\vect{B})$ 项相互抵消：
\[
\frac{\partial}{\partial t}\left(\frac{1}{2}\rho_m u^2 + \frac{B^2}{2\mu_0}\right) = -\divg\left(\frac{1}{2}\rho_m u^2\vect{u} + \vect{S}\right) - \vect{u}\cdot\grad p
\]

$\vect{u}\cdot(\vect{j}\times\vect{B})$ 是磁能与动能之间的\textbf{转换项}（内部交换，不改变两者之和）：
\begin{itemize}[leftmargin=*,itemsep=0pt]
\item 当 $\vect{u}\cdot(\vect{j}\times\vect{B}) > 0$：洛伦兹力对等离子体做正功，\textbf{磁能转化为动能}（如磁驱动喷流、等离子体加速）
\item 当 $\vect{u}\cdot(\vect{j}\times\vect{B}) < 0$：等离子体做功抵抗磁场，\textbf{动能转化为磁能}（如发电机效应、磁场放大）
\end{itemize}

\textbf{但这还不是完整的能量守恒}——上式右端仍留着 $-\vect{u}\cdot\grad p$，它把能量交给了\textbf{内能}，而内能没有出现在守恒式里。磁能与动能之和并不单独守恒：压强功会把能量泵入内能，或反过来。\textbf{加上内能才闭合}。

\textbf{提高题（完整能量守恒）：} 对绝热理想MHD，内能密度为 $U = p/(\gamma-1)$，其演化方程为
\[
\frac{\partial U}{\partial t} + \divg(U\vect{u}) = -p\divg\vect{u}
\]
把三项相加，$-\vect{u}\cdot\grad p$ 与 $-p\divg\vect{u}$ 正好组合成 $-\divg(p\vect{u})$，于是得到\textbf{封闭形式}：
\[
\boxed{\;\frac{\partial}{\partial t}\left(K + U + U_B\right) + \divg\Bigl\{\vect{u}\left(K + U + p\right) + \vect{S}\Bigr\} = 0\;}
\]
其中 $K=\tfrac12\rho_mu^2$、$U=p/(\gamma-1)$、$U_B=B^2/(2\mu_0)$。请自行验证：$\vect{u}(K+U+p)$ 是随流体带走的对流焓流，$\vect{S}$ 是坡印廷能流。

\textbf{三层含义要分清：}
\begin{itemize}[leftmargin=*,itemsep=0pt]
\item \textbf{局域交换}：$\vect{u}\cdot(\vect{j}\times\vect{B})$ 与 $-\vect{u}\cdot\grad p$ 是同一空间点上不同能量库之间的转移，不进出系统；
\item \textbf{边界能流}：散度项 $\divg\{\cdots\}$ 描述能量跨越控制面的输运；
\item \textbf{封闭系统总量}：对全空间积分，散度项化为边界通量。取边界上无能量流出，则 $\frac{\diff}{\diff t}\int(K+U+U_B)\diff V=0$——\textbf{这才是"能量守恒"四个字所指的结论}。
\end{itemize}
只把 $K$ 与 $U_B$ 相加就宣称"总能量守恒"，遗漏了内能与之相伴的压强功，是初学者最常见的误读。这一条也直接连到第10章的加热与第11章的输运损失：能量进了内能，究竟是被带走还是被辐射掉，正是那里要回答的问题。
\end{example}

\begin{figure}[htbp]
\centering
\begin{tikzpicture}[scale=1.0]
% 左侧：流体元和磁力线
\draw[thick,fill=orange!10] (-2,0) rectangle (0,1.5);
\node at (-1,0.75) {\small 流体元};

% 磁力线穿过流体元
\foreach \y in {0.3,0.75,1.2} {
  \draw[->,blue,thick] (-3,\y) -- (-2,\y);
  \draw[blue,thick] (-2,\y) -- (0,\y);
  \draw[->,blue,thick] (0,\y) -- (1,\y);
}
\node[blue] at (1.4,0.75) {$\vect{B}$};

% 速度箭头
\draw[->,red,very thick] (-1,1.8) -- (1,1.8);
\node[red,above] at (0,1.9) {$\vect{u}$};

% 右侧：移动后的流体元（虚线）
\draw[thick,fill=orange!10,dashed] (0.5,0.3) rectangle (2.5,1.8);
\node[orange!80!black] at (1.5,1.05) {\small 流体元};

% 移动后的磁力线（虚线）
\foreach \y in {0.6,1.05,1.5} {
  \draw[->,blue,thick,dashed] (0.5,\y) -- (2.5,\y);
  \draw[->,blue,thick] (2.5,\y) -- (3.5,\y);
}
\node[blue,dashed] at (3.9,1.05) {$\vect{B}$};

% 标注
\draw[<->,gray] (-2,-0.3) -- (0.5,-0.3);
\node[gray] at (-0.75,-0.6) {$\vect{u}\,\diff t$};

\node at (0,-1.2) {\small 理想MHD：流体元带着磁力线一起运动，穿过任意回路的磁通量不变};
\end{tikzpicture}
\caption{磁力线冻结示意图。在理想导电流体（$R_m\gg 1$）中，磁通量随流体元守恒，磁力线仿佛被"冻结"在流体中。}
\end{figure}

\begin{figure}[htbp]
\centering
\begin{tikzpicture}[scale=1.0]
% 弯曲磁力线（向上凸起）
\draw[blue,very thick,domain=-3:3,samples=100] plot(\x,{0.5*cos(30*\x)});
\draw[blue,->,very thick] (-2,0.43) -- (-1.9,0.42);
\draw[blue,->,very thick] (0,0.5) -- (0.1,0.49);
\draw[blue,->,very thick] (2,0.43) -- (2.1,0.42);
\node[blue] at (3.2,0.5) {$\vect{B}$};

% 曲率中心
\fill[black] (0,6) circle (2pt);
\node at (0,6.5) {曲率中心};
\draw[dashed] (0,6) -- (-1,0.43);
\draw[dashed] (0,6) -- (1,0.43);

% 张力恢复力（指向曲率中心）
\draw[->,red,very thick] (-1,0.43) -- (-0.8,1.5);
\node[red] at (-0.5,1.8) {$\vect{F}_\text{张力}$};
\draw[->,red,very thick] (1,0.43) -- (0.8,1.5);
\node[red] at (0.5,1.8) {$\vect{F}_\text{张力}$};

% 磁压方向（向下，推向B弱区）
\draw[->,green!60!black,very thick] (-1.5,2.5) -- (-1.5,1.5);
\draw[->,green!60!black,very thick] (1.5,2.5) -- (1.5,1.5);
\node[green!60!black] at (2.2,2) {$\grad(B^2/2\mu_0)$};

% 等离子体区域
\fill[orange!10] (-3,-1) rectangle (3,-0.2);
\node at (0,-0.6) {等离子体};

\node at (0,-1.8) {\small 磁张力像拉紧的橡皮筋，将弯曲磁力线拉直；磁压将等离子体推离高场区};
\end{tikzpicture}
\caption{磁压与磁张力示意图。弯曲磁力线产生指向曲率中心的张力恢复力，磁压梯度将等离子体推向磁场较弱的区域。}
\end{figure}

% ch5 新例题：Z-pinch 平衡（Bennett 压缩）
\begin{example}{Z-pinch 的 Bennett 关系与平衡温度}{bennett-pinch}
一直圆柱等离子体（Z-pinch）流过轴向电流 $I = 100\,\mathrm{kA}$，线密度（每单位长度粒子数）为 $N = 10^{19}\,\mathrm{m^{-1}}$，离子为氘（$m_i = 2m_p$）。用 Bennett 关系估算维持平衡所需的温度 $T_e + T_i$，并讨论该平衡为何不稳定。
\tcblower
\textbf{解答：}

\textbf{Bennett 关系：} Z-pinch 中轴向电流产生角向磁场 $B_\theta(r) = \mu_0 I(r)/(2\pi r)$，磁压 $B_\theta^2/(2\mu_0)$ 向内压缩等离子体（pinch 效应），等离子体热压必须与之平衡。对半径积分后，平衡条件为 \textbf{Bennett 关系}：
\[
\mu_0 I^2 = 8\pi N k_{\!B}(T_e + T_i)
\]
（$N = \int n\,2\pi r\,\diff r$ 为线密度。）

代入数值（$I = 10^5\,\mathrm{A}$，$N = 10^{19}\,\mathrm{m^{-1}}$）：
\[
k_{\!B}(T_e+T_i) = \frac{\mu_0 I^2}{8\pi N}
= \frac{4\pi\times10^{-7}\times10^{10}}{8\pi\times10^{19}}
= \frac{4\pi}{8\pi}\times10^{-7+10-19}
= \frac{1}{2}\times10^{-16}
= 5.0\times10^{-17}\,\mathrm{J}
\]
\[
T_e + T_i = \frac{5.0\times10^{-17}}{1.381\times10^{-23}}\,\mathrm{K} = 3.6\times10^{6}\,\mathrm{K} \approx 310\,\mathrm{eV}
\]
（即 $k_{\!B}T \approx 310\,\mathrm{eV}\times1.602\times10^{-19} = 5.0\times10^{-17}\,\mathrm{J}$。）

即平衡要求 $T_e + T_i \approx 310\,\mathrm{eV}$（电子与离子各约 $155\,\mathrm{eV}$）。这与实际 Z-pinch 实验的温度量级（$100\sim1000\,\mathrm{eV}$）一致，说明 Bennett 关系对 Z-pinch 参数的估计是合理的。反过来，若给定温度 $T\sim100\,\mathrm{eV}$，维持平衡所需的最小电流为 $I \approx 57\,\mathrm{kA}$。

\textbf{不稳定性：} Z-pinch 平衡对 $m=0$（sausage，腊肠模）和 $m=1$（kink，扭曲模）扰动都不稳定（见第9章）：电流产生的角向磁场在截面收窄处增强（磁压更大），进一步压缩——正反馈。稳定 Z-pinch 需要纵向磁场（shear 场）、导电壁或快速等离子体加热。这也是 Z-pinch 至今未成为主流聚变约束方案的核心原因。
\end{example}

\begin{formulabox}{第5章核心公式速查}{ch5-formulas}
\begin{itemize}[leftmargin=*,itemsep=0pt]
\item 理想MHD质量守恒：$\displaystyle\frac{\partial\rho_m}{\partial t} + \divg(\rho_m\vect{u}) = 0$
\item 理想MHD动量方程：$\displaystyle\rho_m\left(\frac{\partial\vect{u}}{\partial t} + \vect{u}\cdot\grad\vect{u}\right) = -\grad p + \vect{j}\times\vect{B}$
\item 理想欧姆定律：$\displaystyle\vect{E} + \vect{u}\times\vect{B} = 0$
\item 磁场演化：$\displaystyle\frac{\partial\vect{B}}{\partial t} = \curl(\vect{u}\times\vect{B})$
\item 磁压：$\displaystyle p_B = \frac{B^2}{2\mu_0}$，磁张力：$\displaystyle\frac{1}{\mu_0}(\vect{B}\cdot\grad)\vect{B}$
\item 比压：$\displaystyle\beta = \frac{2\mu_0 p}{B^2}$
\item 阿尔芬速度：$\displaystyle v_A = \frac{B}{\sqrt{\mu_0\rho_m}}$
\item 阿尔芬波色散关系：$\displaystyle\omega = k_\parallel v_A$（无色散；群速度 $\vect{v}_g=\pm v_A\unitvec{B}$ 恒沿磁场，相速度 $v_p=v_A\cos\theta$ 仅在 $\vect{k}\parallel\vect{B}$ 时等于 $v_A$）
\item 磁雷诺数：$\displaystyle R_m = \frac{\mu_0 v L}{\eta} = \frac{\tau_\eta}{\tau_{\text{流}}}$
\item 静态平衡：$\displaystyle\grad p = \vect{j}\times\vect{B}$（Grad-Shafranov方程的基础）
\end{itemize}
\end{formulabox}

\begin{checklistbox}{第5章自学检查清单}{ch5-checklist}
\begin{itemize}[leftmargin=*,label=$\square$]
\item 我能从双流体方程出发，说明MHD近似的适用条件（$\tau\gg 1/\omega_{ci}$，$L\gg\rho_{Li}$）
\item 我能计算比压 $\beta$、阿尔芬速度 $v_A$ 和磁雷诺数 $R_m$，并判断其物理意义
\item 我理解了磁压与磁张力的分解：洛伦兹力 = 磁压梯度 + 磁张力
\item 我能从$R_m$的大小判断磁力线冻结是否成立，并知道电阻会破坏冻结
\item 我掌握了阿尔芬波的物理图像：磁张力作为恢复力的横波，能量沿磁力线以 $v_A$ 输运（而非沿波矢）
\item 我能推导理想MHD中磁能与动能转换的方程，并解释$\vect{u}\cdot(\vect{j}\times\vect{B})$的符号含义
\item 我知道静态平衡$\grad p = \vect{j}\times\vect{B}$是托卡马克和仿星器设计的出发点
\end{itemize}
\end{checklistbox}

\endinput
